% Оформление печатного варианта ОГЭ. Воспроизводит эталон ФИПИ: тот же
% pdfTeX, те же шрифты cm-super, те же рамки и линейки. Проверка совпадения —
% scripts/print/compare.py (см. docs/print-variants.md).
\usepackage[T2A]{fontenc}
\usepackage[utf8]{inputenc}
\usepackage[russian]{babel}
\usepackage{amsmath,amssymb}
\usepackage{graphicx}
\usepackage{tikz}
\usepackage[most]{tcolorbox}
\usepackage{fancyhdr}
\usepackage{array}
\usepackage{multirow}
\usepackage{wrapfig}
\usetikzlibrary{calc,arrows}

\geometry{
  a4paper,
  left=2cm, right=2cm,
  top=\OGEtop, bottom=\OGEbottom,
  headheight=\OGEheadheight, headsep=\OGEheadsep,
  footskip=\OGEfootskip
}

% Эталон набран с поднятым порогом: есть строки с межсловным пробелом 6,6bp
% при норме 3,9. На умолчании 200 TeX такие строки отвергает и рвёт абзацы
% иначе — совпадения по строкам не будет.
\tolerance=9999

% --- колонтитул ---
\pagestyle{fancy}
\fancyhf{}
\fancyhead[L]{\OGEheadleft}
\fancyhead[R]{\OGEheadright}
\fancyfoot[C]{\thepage}
% Полосы не растягиваются до низа: в эталоне ФИПИ последняя страница
% заканчивается там, где кончились задания, а не растянута на весь лист.
\raggedbottom
\renewcommand{\headrulewidth}{0.4pt}
\renewcommand{\footrulewidth}{0pt}

% Внутри задания разрывы страницы запрещены, во вводном условии — нет:
% условие бывает почти в полосу высотой, и запрет там дал бы переполненную
% полосу вместо аккуратного переноса.
\newif\ifogebind
\newcommand{\ogebind}{\ifogebind\nopagebreak\fi}

% --- плашка с инструкцией: рамка 1pt, фон yellow!1 ---
\newtcolorbox{instrbox}{
  colback=yellow!1, colframe=black,
  boxrule=1pt, arc=0pt, outer arc=1pt,
  left=0.5cm, right=0.5cm, top=\OGEinstrtop, bottom=\OGEinstrbot,
  boxsep=0pt, width=\textwidth,
  before skip=\OGEinstrbefore, after skip=\OGEinstrafter
}

% --- номер задания: внешний габарит ровно 1cm x 19.708bp, линии 1pt ---
% Рисуется четырьмя отрезками, а не rectangle: у прямоугольника обводка идёт
% по центру пути и габарит выходит на полтолщины шире с каждой стороны.
\newcommand{\znum}[1]{%
  \begin{tikzpicture}[baseline=(zz.base),x=1bp,y=1bp]
    \useasboundingbox (0,-0.498) rectangle (28.346,19.209);
    \draw[line width=1pt] (0,0) -- (28.346,0);
    \draw[line width=1pt] (0,18.711) -- (28.346,18.711);
    \draw[line width=1pt] (0.498,0) -- (0.498,18.711);
    \draw[line width=1pt] (27.848,0) -- (27.848,18.711);
    \node[anchor=base,inner sep=0pt] (zz) at (14.173,5.48) {\bfseries #1};
  \end{tikzpicture}%
}

% --- задание: номер в рамке слева, тело с отступом 2cm ---
% raisebox с нулевой глубиной обязателен: иначе рамка добавляет 5,48bp
% глубины первой строке, и строки с «й»/«б» уезжают вниз на полтора пункта.
\newlength{\zbody}\setlength{\zbody}{2cm}
\newenvironment{zadanie}[2][0pt]{%
  \par\penalty-150\vspace{\OGEzskip}\vspace{#1}%
  \noindent
  \begin{list}{}{%
    \setlength{\leftmargin}{\zbody}%
    \setlength{\labelwidth}{\zbody}%
    \setlength{\labelsep}{0pt}%
    \setlength{\itemindent}{0pt}\setlength{\listparindent}{0pt}%
    \setlength{\topsep}{0pt}\setlength{\parsep}{0pt}\setlength{\partopsep}{0pt}%
    \setlength{\itemsep}{0pt}\setlength{\rightmargin}{0pt}}%
  \item[{\raisebox{0pt}[\height][0pt]{\makebox[\zbody][l]{\znum{#2}}}}]%
  % Задание не рвётся между страницами. Разорванное задание — худшее, что
  % может случиться с печатной работой: половина условия на одном листе,
  % чертёж и строка ответа на другом. Самое высокое задание банка (условие
  % плюс чертёж на 240pt) занимает меньше половины полосы, поэтому целиком
  % переносится всегда. Между заданиями разрыв остаётся свободным.
  \interlinepenalty=10000\relax\ogebindtrue
}{\end{list}}

% --- строка «Ответ: ____»: линейка 2.75cm, 0.4pt, серая 50% ---
% Строка ответа тоже привязана к заданию: оторванное «Ответ: ____» в начале
% страницы не подскажет ученику, к чему оно относится.
\newcommand{\otvet}[1][0pt]{%
  \par\nopagebreak\vspace{\OGEotvetskip}\vspace{#1}%
  \noindent\hspace*{\fill}Ответ:\kern2.709bp%
  \textcolor[gray]{0.5}{\rule[-0.76bp]{77.954bp}{0.398bp}}\kern0.202bp\par}

% --- заголовок работы и шапка части ---
\newcommand{\ogeTitle}[1]{%
  \begin{center}
    {\Large\bfseries #1}\\[11.862pt]
    {\bfseries Часть \textnumero{} 1}
  \end{center}}
\newcommand{\ogePartTwo}{%
  \begin{center}{\bfseries Часть \textnumero{} 2}\end{center}}

% --- иллюстрация задания ---
% \nopagebreak с обеих сторон обязателен: между условием и его чертежом
% пролегает законный разрыв страницы, и TeX им пользуется — рисунок к десятому
% заданию уезжает на следующий лист, а само задание остаётся на предыдущем.
\newcommand{\ogefig}[2]{%
  \par\nopagebreak\vspace{\OGEfigbefore}%
  \noindent\hspace*{\fill}\includegraphics[width=#2]{#1}\hspace*{\fill}%
  \par\nopagebreak\vspace{\OGEfigafter}}

% Чертёж задания, свёрстанный отдельным блоком: держится за своё условие.
\newcommand{\ogeTaskFig}[2]{% {файл}{ширина pt}
  \par\nopagebreak\vspace{\OGEfigbefore}%
  \noindent\hspace*{\fill}\includegraphics[width=#2pt]{#1}\hspace*{\fill}%
  \par\nopagebreak\vspace{\OGEfigafter}}

% Чертёж справа от текста. Ширина текстовой колонки — как в эталоне ФИПИ:
% текст занимает 311.8bp, остальное отдано чертежу.
\newlength{\ogetextcol}\setlength{\ogetextcol}{311.8bp}
\newcommand{\ogeSideBySide}[2]{%
  \par\ogebind\noindent
  \begin{minipage}[t]{\ogetextcol}#1\end{minipage}\hfill
  \begin{minipage}[t]{\dimexpr\linewidth-\ogetextcol-8pt\relax}%
    \vspace{0pt}\centering #2%
  \end{minipage}\par}

% Таблица данных из условия — по центру тела задания.
\newcommand{\ogeTable}[2]{%
  \par\ogebind\vspace{\OGEtabbefore}%
  \noindent\hspace*{\fill}\begin{tabular}{#1}#2\end{tabular}\hspace*{\fill}%
  \par\ogebind\vspace{\OGEtabafter}}

% Раскладочная таблица ФИПИ — без линеек и с узкими промежутками: ею
% свёрстаны полосы «А) чертёж Б) чертёж» и подписи под рисунками.
\newcommand{\ogePlainTable}[2]{%
  \par\ogebind\vspace{4pt}%
  \noindent\hspace*{\fill}%
  {\setlength{\tabcolsep}{3pt}\begin{tabular}{#1}#2\end{tabular}}%
  \hspace*{\fill}\par\ogebind\vspace{4pt}}

% Варианты ответа к заданию с выбором.
\newcommand{\ogeChoices}[1]{\par\ogebind\vspace{4pt}#1}
\newcommand{\ogeChoice}[2]{%
  \par\ogebind\noindent\makebox[16pt][l]{#1)}%
  \begin{minipage}[t]{\dimexpr\linewidth-16pt\relax}#2\end{minipage}\par\ogebind}

% --- таблицы данных ---
\setlength{\arrayrulewidth}{0.4pt}
\newcommand{\cs}{\rule[-7.0pt]{0pt}{23.036pt}}

% Вводное условие блока 1–5: без рамки и номера, во всю ширину полосы.
\newcommand{\ogeIntro}[1]{\par\vspace{4pt}#1\par\vspace{6pt}}

% Растровая иллюстрация условия (планы, чертежи, фотографии ФИПИ).
% Печатается в натуральную величину и только сжимается, если не влезает в
% полосу: растянутый растр читается хуже мелкого, а увеличивать нечего —
% исходник у ФИПИ экранный.
\newcommand{\ogeRaster}[2]{%
  \par\vspace{\OGEfigbefore}%
  \noindent\hspace*{\fill}%
  \ifdim #2pt>\linewidth
    \includegraphics[width=\linewidth,height=0.33\textheight,keepaspectratio]{#1}%
  \else
    \includegraphics[width=#2pt,height=0.33\textheight,keepaspectratio]{#1}%
  \fi
  \hspace*{\fill}\par\vspace{\OGEfigafter}}

% Та же иллюстрация внутри ячейки таблицы: без отбивок и переводов строки,
% иначе \par порвёт ячейку.
\newcommand{\ogeRasterCell}[2]{%
  \ifdim #2pt>0.42\linewidth
    \includegraphics[width=0.42\linewidth,height=0.22\textheight,keepaspectratio]{#1}%
  \else
    \includegraphics[width=#2pt,height=0.22\textheight,keepaspectratio]{#1}%
  \fi}

% Иллюстрация условия с обтеканием: рисунок справа, текст идёт слева от него.
% Так свёрстаны практико-ориентированные блоки в печатных сборниках — иначе
% рисунок 200pt шириной забирает всю полосу и работа распухает на страницу.
% Первый параметр — сколько строк текста обтекает: wrapfig не умеет считать
% это сам и на глазок оставляет либо дыру, либо наползание.
\newcommand{\ogeWrap}[4]{% {строк}{ширина pt}{файл}{подпись}
  \begin{wrapfigure}[#1]{r}{\dimexpr #2pt+10pt\relax}
    \vspace{-\intextsep}
    \centering
    \includegraphics[width=#2pt]{#3}%
    \ifx\relax#4\relax\else\par\vspace{2pt}{\footnotesize #4}\fi
  \end{wrapfigure}}

% Обозначение внутри предложения: «имеет ширину [B=195] мм». Такие формулы
% ФИПИ тоже хранит растром, и вынести их отдельной иллюстрацией — значит
% выбросить кусок предложения и оставить «имеет ширину мм».
\newcommand{\ogeRasterInline}[2]{%
  \raisebox{-0.25\height}{\includegraphics[height=#2pt]{#1}}}

% Место под решение во второй части.
\newcommand{\ogesolutiongap}[1][\OGEparttwogap]{\par\vspace{#1}}

% Слова, которые эталон не переносит: его генератор пользовался другим
% набором русских паттернов, и без этих исключений абзацы рвутся иначе.
\hyphenation{сколько от-но-си-тельно}
